\chapter*{Preface}
\addcontentsline{toc}{chapter}{Preface}
MoDELib---the Mechanics of Defects Evolution Library---provides a framework for modeling defects in crystalline materials and their interactions with mechanical and other physical fields. Understanding such simulations requires both a physical description of the problem and a clear account of how that description is represented in the code. My aim with this manual is to bring these two perspectives together.

I have organized the manual around the architecture of the library. It begins with the geometric foundations: simplicial and polyhedral meshes, their connectivity, and periodic identification. Materials and polycrystals then give this geometry its crystallographic meaning. The finite-element infrastructure supplies interpolation, weak forms, and field constraints, while the multiphysics architecture connects the individual physical models through shared data and a common simulation lifecycle.

This order is intended to help readers build an understanding of the software before entering the details of a particular model. Theory is introduced where it explains an implementation choice or a physical assumption. Figures and short code examples connect the mathematical objects to their computational counterparts. The appendices collect longer dislocation-theory derivations and supporting identities, so that these remain available without interrupting the architectural discussion.

The manual is also a working reference. Some interfaces continue to evolve, and the cluster-dynamics formulation is being revised. The code version identified on the cover therefore provides the context for the descriptions and examples presented here. As the library develops, I intend this document to preserve the connection between the governing equations, the design of the software, and the practical choices needed to construct and interpret a simulation.

\vspace{1.2cm}
\begin{flushright}
\textit{Giacomo Po}\\[0.3em]
21 October 2026
\end{flushright}
